(1) 电子与核自旋均为 \(1/2\)，哈密顿量为 \(H=A\mathbf S\cdot\mathbf I/\hbar^2+b_s S_z/\hbar+b_i I_z/\hbar\)，其中 \(A>0\)、\(b_s=\alpha_sB,b_i=\alpha_iB\)。按 \(|\uparrow\uparrow\rangle,|\uparrow\downarrow\rangle,|\downarrow\uparrow\rangle,|\downarrow\downarrow\rangle\) 顺序建立 \(4\times4\) 矩阵及分块。
(2) 完整求出四支能谱和 \(m=0\) 的本征态投影算符；讨论 \(B=0\) 单重/三重态间距及 \(|b_s-b_i|\gg A\) 极限。
(3) 弱横向探针只通过电子 \(S_x\) 跃迁。求初态 \(|\uparrow\uparrow\rangle\) 到两个 \(m=0\) 本征态的强度、和则及零场与强场的选择规则。
(4) 求任意 \(T>0\) 的配分函数、磁化强度与零场磁化率，证明磁化率非负，求高温Curie律和低温Van Vleck极限。
